\section{Installing the original quotient slopes and frame inequality}
\label{sec:global-assembly-v034}

\status{The following arguments install the finite-place and frame inputs
on the original supplied family.  They keep the arithmetic inputs
\texttt{LinkG}, \texttt{LinkS}, \texttt{LinkIRR}, \texttt{LinkDIV} and
\texttt{LinkMu} as inputs.  In particular, they do not manufacture the
common second layer or the spectral estimates.  The remaining final
assembly is specified at the end of the chapter.}

Fix an original family $\mathcal F:\texttt{FamP}(\epsilon)$,
$0<\epsilon<1/52$, and its fixed constants.  A statement asserted
eventually means that a single threshold is chosen before the index
$i\in\mathcal F.I_x$, the word $u$, and every finite place.  We write
$D=\mathcal F.D(x,i)$ and $L=\mathcal F.L(x,i)$ only after that choice.
The scales retain their original meanings:
\[
 b=e^{x^\epsilon},\qquad a=xb,\qquad
 X'=e^{x^{\epsilon/3}},\qquad b'=a/X',\qquad
 X''=xnX'.                                                    \tag{GFA.1}
\]
For a word let $M_u$ be the original ramified reference lattice,
$N_u^0=L.\texttt{NuLat}(u)$, and
$N_u=D.\texttt{div.Nu}(u)=I_uN_u^0$, with the unchanged metrics of MT1.
Write $d_u=D.\texttt{duS}(u)$, $r_u=\operatorname{rank}N_u$ and
$T=T_u$.  The full source lattice is denoted $A_u=(D.\texttt{fm}(u)).
\texttt{lat}(0)$.  Thus $A_u$ and $M_u$ are not confused.

% BEGIN GFA1
\paragraph{GFA1: the exact type-prime cost, with constant $36$.}
\label{claim:GFA1}
Assume the original ramified-model identities and congruence-length
bound supplied by RSL9--RSL10 and GLE11, together with
\texttt{TuSize}.  GLE2 activates the guarded statements simultaneously.
The full numerical and weight windows from GVC11's
\texttt{NumsL}/\texttt{WtsL} installation may be taken with $c_L=3$.
Here the needed logarithmic-derivative bound is already the
\texttt{DlSize} clause of the original size record.

We first check the numerical factor in the congruence estimate.  The
actual top-coordinate identity gives $a_{\max}\le2a$.  Every other
entry of the lifted weight at a top has absolute value at most $2b'$:
the interior and last entries have this bound by \texttt{WtsL}; the
absolute value of each appended block entry lies between the
bottom weight and the largest non-thick $w_j\le2b'$.  The determinant shift
$i(t+1)/2$ is at most $N(m)\le b'$.  Consequently the original
shifted absolute weight $\operatorname{Wt}$ is at most $3b'$.
Also $k\le N(m)\le b'$ and $b_{\rm top}\le2b'$.
At any selected place over $T$, its summand therefore satisfies
\[
 \frac{a_{\max}}k+2\operatorname{Wt}+k+b_{\rm top}
 \le 2a/k+9b'\le9(a/k+b').                              \tag{GFA.2}
\]
The map from the selected local positions to the original tops is
injective (it is the restriction of the actual $\tau$-bijection).
There are at most $|TT|$ such positions.  Hence (GFA.2) proves
\texttt{LevelPrimeSizeBounds} with $C_1=9$ on the original data.

Let $\ell_w=\log N\mathfrak p_w/[K:\mathbb Q]$ and
$e_w=v_w(T)$ for $w\mid T$, where $K=D.\texttt{what}_2.LL$.
For a finite set $F$ of such places put
$B=C_1|TT|(a/k+b')$.  GLE11 bounds the actual local quotient length by
\[
 n_w\le e_w m d_u B.
\]
The source's \texttt{NuDef} and \texttt{LenAtT} identify this same
integer with $e_wL.\texttt{lenMN}(u,w)$; no alternative quotient is
substituted.  The degree identity of LC3 now gives
\[
 \sum_{w\in F}2\ell_we_wL.\texttt{lenMN}(u,w)
 \le2md_uB\sum_{w\in F}e_w\ell_w
 \le2md_uB\log T.                                      \tag{GFA.3}
\]
The setting supplies $T>2fN(m)+1\ge2f+1$.  By Bernoulli,
\[
 q_u:=\left(1-\frac1{T-1}\right)^f
 \ge1-\frac f{T-1}\ge\frac12.
\]
There is a rank identification to check before this division.
GLE11 gives the original inclusion $N_\lambda\subset S_{\rm ram}$.
On this same supplied family it activates RSL11, hence fullness
of $N_u^0$ and the actual \texttt{NuFull} property of $N_u$.
The source identities
\texttt{NuV\_finrank} and \texttt{NuFull.rank} then identify the
rank of the actual $N_u$ with the ambient retained dimension
$(D.\texttt{mem}(u)).A.\texttt{du}$.  RSL10 computes that dimension
as $d_uq_u$.  Thus $r_u=d_uq_u$ for the very lattice used here.
Dividing (GFA.3) by $r_u$ yields
\[
 \boxed{\quad
 \frac1{r_u}\sum_{w\in F}2\ell_we_wL.\texttt{lenMN}(u,w)
 \le36|TT|m(a/k+b')\log T.
 \quad}                                                 \tag{GFA.4}
\]
If $d_u=0$, the numerator and retained rank vanish; the formal
total-division convention gives zero on the left.  Every later use
has $r_u>0$, so uses the ordinary division just written.

The other four clauses of the original \texttt{adm26} are also
accounted for.  Its rank, mean and maximal-slope clauses are exactly
RSL10.  For the local inclusion at $T$, \texttt{NuAtTDef} maps
$(N_u^0)_w$ into the same congruence lattice $N_\lambda$.
GLE11 gives $N_\lambda\subset S_0$; the original localized class
lattice identity pulls this back to $(M_u)_w$.  Together with
(GFA.4), this proves \texttt{LevelTFamily} with constant $36$.
Taking the maximum of the setting, numerical, weight, and
congruence thresholds occurs before $i,u,w$.
% END GFA1

% BEGIN GFA2
\paragraph{GFA2: the required total correction bound directly from its definitions.}
\label{claim:GFA2}
The quotient-slope argument needs only
\[
 \mathrm{Ta}:=\sum_{w\in L.\mathrm{supp}}2\ell_we_w\tau_w\le b^2.
                                                               \tag{GFA.5}
\]
We prove this on the actual family without assuming the stronger
\texttt{FamilyWeightSumBound}.  This is a deliberately coarse
estimate for this particular consumer.

Set $S=\exp(x^{\epsilon/2})$.  The original rank windows, GII2--GII3,
and the elementary comparison of fixed powers with exponentials give,
eventually and uniformly,
\[
\begin{gathered}
 S\ge40,\quad a\ge1,\quad
 f,m,n,k,N_1,r_G,x^{26},X',\log a/\log2\le S,\\
 \max(1,|C_e|,\text{each fixed local constant})\le S,
 \qquad 0\le c_G=4L.\mathrm{Sc}\le4b'/x\le5a,\quad b'\le a,\\
 \Lambda:=C_{\rm La}N_1^2\log a\le S^4,
 \qquad 0\le c_0=2b'+2c_G\le12a .
\end{gathered}                                               \tag{GFA.6}
\]
For every $p<X'$, GII2 gives
$G_p\le C_e x^{\epsilon/8}\log x$.  Therefore
\[
 \log(p^{G_p})\le C_e x^{11\epsilon/24}\log x
 \le x^{\epsilon/2},\qquad Q_p=p^{G_p}\le S.                \tag{GFA.7}
\]
The strict exponent difference is $\epsilon/24>0$; the threshold
depends only on the fixed $C_e,\epsilon$.  This is also valid for
the original natural ceilings, already included in GII2.

If the original set $VV(p)$ is nonempty, its definition implies that
$p$ is prime, $p<X'$, and it is inert.  The exact orbit identification
then gives $n_p=\operatorname{ord}(h_p)\le n\le S$.
There are at most $n$ representatives and at most $n$ tops in any
orbit.  Since $\delta_1\le1$,
$N^+=40n_p\delta_1\le S^2$.  We have
$\log_p a\le S$, $p^{-L_1}\le1$, $pk\ge1$, and
$0\le b_t\le2b'\le2a$ at every top.  These are the only
orbit and weight facts used below; no new good-orbit estimate is
assumed.

Here are upper bounds for every term in the literal definitions.
They keep the two branches of $D_p,t_v$ distinct:
\[
\begin{array}{c|c|l}
\text{quantity}&\text{upper bound}&\text{reason before rounding}\\\hline
 D^{\rm I}_p&S^{11}&
 C_{DI}N^+x^{26}(N_1^5+k^3)\log_pa\le2S^{10}\\
 D^{\rm II}_p&aS^{11}&
 3S^9+10aS^2+12aS^3\\
 t^{\rm II}_v&aS^8&
 S^3(S^4a+a+12a+S^2)+10aS^2\\
 c_6^{\rm I}&aS^{12}&
 8(S^{11}+1)+S^5+20aS^2+10aS\\
 t^{\rm I}_v&aS^{16}&
 2aS+2aS^7+2aS^{15}+10aS .
\end{array}                                                  \tag{GFA.8}
\]
For the first two rows use $\lceil z\rceil_+\le\max(z,0)+1$.
The inequalities displayed follow for $S\ge40,a\ge1$.
For the last row, the predecessor sum has at most $S$ terms and
each term is bounded by its bottom weight $2a$; if its distance
is zero the term is zero.  Its second summand is at most
$S^3(a+S^4a)$.  The third is at most
$S^2(S\cdot aS^{12}+2S^2+5a)$.  This proves the final row,
including the actual minimum and distance convention.

The stored ramification exponent is
$e_p=\lceil C_{\rm ram} f k^2\log_pa\rceil_+$, where
$C_{\rm ram}$ is the fixed original \texttt{cX.Ce} (distinct
from the orbit constant $C_e$).  Thus $e_p\le2S^5$.
The original definition
\[
 \tau_p=e_p+\sum_{v\in VV(p)}m(\lceil t_v\rceil_++D_p+1)
\]
and (GFA.8) imply $0\le\tau_p\le aS^{20}$.  When $VV(p)$
is empty the definition instead sets $\tau_p=0$.
In particular every $\tau_p$ is a nonnegative integer.
The original fill equality is $\tau_w=\tau_{p(w)}$.
For the finite support, LC3, nonnegativity, and (GFA.7) now give
\[
 \mathrm{Ta}\le2\sum_{p<X'}\log p\,\tau_p
 \le4aS^{22}\le aS^{24}.
\]
Indeed there are at most $\lceil X'\rceil\le2S$ eligible primes,
and $\log p\le S$.  Finally
\[
 \log(aS^{24})=x^\epsilon+\log x+24x^{\epsilon/2}
 \le2x^\epsilon=\log b^2
\]
eventually.  This proves (GFA.5) with one threshold for every
original member and word.  It neither proves nor needs the
sharper Mertens-weighted estimate of the alternative source route.
% END GFA2

% BEGIN GFA3
\paragraph{GFA3: exact mean and maximum slopes from all finite places.}
\label{claim:GFA3}
Let $M$ be a full positive-rank hermitian lattice over $K$, and
$N^0\subset M$ a commensurable lattice with the same metric.
For $n_w=\operatorname{length}_{\mathcal O_w}(M_w/N_w^0)$,
local elementary divisors and the norm of the finite quotient give
\[
 \mu(N^0)=\mu(M)-\frac1r\sum_w2\ell_wn_w.                  \tag{GFA.9}
\]
Only finitely many $n_w$ are nonzero: a common nonzero denominator
contains a multiple of $M$ in $N^0$, and its prime divisors contain
the support.  Multiplication by
$I=\prod_w\mathfrak p_w^{-\lfloor e_ws_w\rfloor}$ gives, by MT2,
\[
 \mu(IN^0)=\mu(M)+\sum_w2\ell_w\lfloor e_ws_w\rfloor
                     -\frac1r\sum_w2\ell_wn_w.              \tag{GFA.10}
\]

For completeness the corresponding maximal-slope inequality is
also local to global.  If $N_w^0\subset\mathfrak p_w^{k_w}M_w$
at a finite set, with ordinary inclusion elsewhere, then
$N^0\subset JM$ for $J=\prod_w\mathfrak p_w^{k_w}$.
To see this, localize the proposed quotient at any maximal ideal;
all other prime-ideal factors become units and the stated local
inclusion applies.  A nonzero quotient element would have a proper
annihilator contained in some maximal ideal, contradicting these
zero localizations.  Taking every nonzero rational subspace, its
induced sublattice in $N^0$ has no larger degree than its induced
sublattice in $JM$.  Taking suprema, and using MT2 on each such
subspace, proves
\[
 \mu_{\max}(IN^0)\le\mu_{\max}(M)
       +\sum_w2\ell_w(\lfloor e_ws_w\rfloor-k_w).             \tag{GFA.11}
\]
This argument also verifies that the scalar-ideal shift for maximum
slope uses the same normalized squared metrics as the mean shift.

On the original records, $\texttt{NuDef}$ reads
$n_w=e_wL.\texttt{lenMN}(u,w)$, and \texttt{NprimeDef} is the
literal equality $N_u=I_uN_u^0$.  RSL11 provides commensurability
and fullness; its proof preserves the ambient $\texttt{Nu.V}$ and
the reference $\texttt{MramL}$.  Thus (GFA.10)--(GFA.11) apply
to these objects, with $r=r_u$, once the support and inclusions
checked next are supplied.
% END GFA3

% BEGIN GFA4
\paragraph{GFA4: the support, integer exponents, and each word's type prime.}
\label{claim:GFA4}
Assume \texttt{LinkG}, \texttt{TuSize}, and the actual family fill,
local-lattice and error-definition fields.  Since $m\to\infty$
in its size window, eventually $m\ge6$ and
\[
 T_u\ge b^{m/2}\ge b^3\ge X'.                              \tag{GFA.12}
\]
The definition of $VV(p)$ forces $p<X'$ when it is nonempty.
Consequently $\tau_{T_u}=0$ for every word, including when
different words have different $T_u$.
GFA2 proves the exact integrality $e_w\tau_w\in\mathbb Z_{\ge0}$,
so $\lfloor e_w\tau_w\rfloor_+=e_w\tau_w$.
The original $\mathrm{io}_p$ is nonnegative: its defining
$\mathrm{Sig}$ is a sum of nonnegative $\kappa$ values.
In the explicit minimum branch this follows from
$\mathcal F.\texttt{GK\_nonneg}$.  In the other branch the
original $\texttt{Extra.kapI\_def}$ makes $\kappa$ the infimum
of admissible nonnegative values.  That set is nonempty:
$\kappa=1$ permits the zero-dimensional choice, whose required
rank bound is $(1-\kappa)\dim=0$ and whose remaining indexed
conditions are empty.  Thus its infimum is also nonnegative.
Furthermore
$S_v\ge0$ because it is the ceiling of a sum of nonnegative
parts.  The empty-set branch is zero.  The fill equality therefore
gives $\mathrm{io}_w\ge0$ and $\mathrm{IOsum}\ge0$.

At $p(w)\ne T_u$, \texttt{LinkG} supplies precisely
\[
 (N_u^0)_w\subset\mathfrak p_w^{e_w\tau_w}(M_u)_w,
 \qquad L.\mathrm{lenMN}(u,w)\le d_u(\tau_w+\mathrm{io}_w).
                                                               \tag{GFA.13}
\]
At $T_u$ use GFA1's inclusion, which has exponent zero, and
retain its separate cost (GFA.4).

It remains to justify deletion of every place outside the stored
finite support.  At a regular such place, support saturation
implies $p(w)>X''\ge X'$, whence $VV(p(w))$ is empty.
The actual \texttt{lats0} field then says
$L.\mathrm{lenMN}(u,w)=0$.  At a type place, use the error field
at the original lift $j=1$ (it exists since $m\ge2$).  If
\texttt{seenP} were true, its $\mathrm{EatT}$ formula would have
the strictly positive summand $fN(m)$ and only nonnegative other
summands.  This contradicts $L.E(u,1,w)=0$ outside support.
Thus \texttt{seenP} is false.  The actual type-place set is then
empty, its exponent is zero and its ordered operator product is
the identity.  Injectivity of the original scalar-extension map
gives $N_\lambda=S_{\rm ram}$.  Its quotient has length zero,
and \texttt{NuAtTDef} identifies it with the required local
length.  This proves the full off-support assertion used in
(GFA.10), at regular and exceptional places alike.
% END GFA4

% BEGIN GFA5
\paragraph{GFA5: the literal \texttt{PlacesF} and \texttt{SlopesNF} inputs.}
\label{claim:GFA5}
Assume, on the same family, \texttt{LevelTFamily},
\texttt{NuCommensF}, \texttt{MuRank}, \texttt{TuSize},
\texttt{LinkG} and \texttt{LinkMu}.  GFA1 supplies the first
with constant $C_{26}=36$ under its stated inputs; alternatively
retain any supplied $C_{26}$.  Write $C_\mu$ for the original
\texttt{LinkMu} constant and choose
$C=\max(C_{26},C_\mu,4)$ before the final threshold.

Put
\[
\begin{gathered}
\delta_u=1-r_u/d_u,\qquad
R_u=\sum_{w\in\mathrm{supp}}2\ell_w
                    (e_ws_{u,w}-\lfloor e_ws_{u,w}\rfloor),\\
B_u=\frac{\sum_{w\in\mathrm{supp},\ p(w)=T_u}
                 2\ell_we_wL.\mathrm{lenMN}(u,w)}{r_u}.
\end{gathered}
                                                               \tag{GFA.14}
\]
Here $r_u>0$, $d_u>0$, $0\le\delta_u<1$ and
$\delta_u\le2f/T_u$ by RSL10 and Bernoulli.  GFA1 gives
$B_u\le C|TT|m(a/k+b')\log T_u$.
The error $R_u$ is nonnegative, and it is exactly the negative of
the signed metric discrepancy in MT3, not merely another upper
bound for it.  Thus \texttt{LinkMu} gives $R_u\le CX''$.

The parameter conditions needed to absorb the retained-rank error
follow directly from (GFA.12).  Eventually $b\ge4$ and $16f\le b$.
Then $2f/T_u\le1/2$, and
\[
 \frac{4fb^2}{T_u}\le\frac{4f}{b}\le\frac14,
 \qquad
 \frac{f\log(T_u-1)}{T_u-2}
 \le\frac{4f}{\sqrt{T_u}}\le\frac14.                       \tag{GFA.15}
\]
For the second inequality use $\log t\le2\sqrt t$ for $t\ge1$,
$T_u-2\ge T_u/2$, and $\sqrt{T_u}\ge b$.
In particular the sum in (GFA.15) is at most $1$, with room to spare.

The exact definitions give
\[
 \sum_w2\ell_we_ws_{u,w}
       =\mathrm{Ga}+\mathrm{Fo}-\mathrm{Ee}_u+\mathrm{Ta}.
\]
At the regular places use (GFA.13), and at the type places
$\tau_w=0$.  Equation (GFA.10) and RSL10 therefore give
\[
\begin{split}
 \mu(N_u)\ge{}&\mu(A_u)+\mathrm{Ga}+\mathrm{Fo}-\mathrm{Ee}_u\\
 &-\frac{2\mathrm{IOsum}}{1-\delta_u}
   -\frac{\delta_u\mathrm{Ta}}{1-\delta_u}
   -\frac{f\log(T_u-1)}{T_u-2}-B_u-R_u\\
 \ge{}&\mu(A_u)+\mathrm{Ga}+\mathrm{Fo}-\mathrm{Ee}_u
             -4\mathrm{IOsum}-B_u-R_u-1 .
\end{split}                                                    \tag{GFA.16}
\]
GFA2 supplies $\mathrm{Ta}\le b^2$ in this calculation.
The factor $2$ before $\mathrm{IOsum}$ is the squared-norm
prime-place normalization, and the second factor $2$ comes from
$1/(1-\delta_u)\le2$.  Both factors are retained.
For the maximum slope take $k_w=e_w\tau_w$ in (GFA.11).
The floor is at most $e_ws_{u,w}$, and RSL10 bounds
$\mu_{\max}(M_u)$ by $\mu_{\max}(A_u)$.  Hence
\[
 \mu_{\max}(N_u)\le
 \mu_{\max}(A_u)+\mathrm{Ga}+\mathrm{Fo}-\mathrm{Ee}_u.
                                                               \tag{GFA.17}
\]

For clarity, the witnesses of the source's existential
\texttt{Places\_at} predicate are
\[
 s'_{u,w}=\frac{\lfloor e_ws_{u,w}\rfloor}{e_w},\quad
 \delta_u,\quad \mu(M_u),\quad\mu_{\max}(M_u),\quad
 P_T=L_T=B_u,\quad R_d=R_X=R_u.
\]
Equations (GFA.10), (GFA.13)--(GFA.16) give its two lower-bound
packages; (GFA.11) gives its upper-bound package.  The last two
inequalities are precisely the stated bounds for $B_u,R_u$.
This checks every witness on the original objects and proves
\texttt{PlacesF}.

With
\[
 \mathrm{Lo}_u=C\mathrm{IOsum}
    +C|TT|m(a/k+b')\log T_u+CX''+1,
\]
(GFA.16) proves the lower clause of \texttt{SlopesN\_at}.
Its upper clause follows from (GFA.17) by adding the nonnegative
allowance $CmB_0$.  This is the literal \texttt{SlopesNF}.
All estimates used finitely many uniform thresholds; their maximum
is chosen before any $i,u,w$.  In particular the shortcut in GFA2
does not change the family or any of these quantifiers.
% END GFA5

% BEGIN GFA6
\paragraph{GFA6: positive target rank and the exact lattice dual.}
\label{claim:GFA6}
Assume \texttt{LinkS}, \texttt{ChosenProviso} with $C_0>0$,
\texttt{LinkIRR} and \texttt{LinkDIV}.  The proviso gives
$d_u\ge D_*/(C_0\Gamma^2)>0$.  The spectral diagonal bound has
\[
 \|K_{uu}-1\|_{\rm HS}^2\le e_0d_u,
 \qquad 0\le e_0,\quad\sqrt{e_0}\le e^{-x^{\epsilon/8}}<1.
\]
If the full analytic target had dimension zero, its period maps
would vanish and $K_{uu}=0$.  The left side would be $d_u$,
contradicting $e_0<1$.  Thus $D_{\rm all}>0$.
The actual target defect obeys
$0\le\delta=D.\texttt{what}_2.\texttt{delp}
 \le e^{-x^{e_z/7}}<1$ by \texttt{LinkIRR}.
The stored equality $D_{\rm rows}=(1-\delta)D_{\rm all}$ and
the fullness of $RR$ from \texttt{LinkDIV} give
$\operatorname{rank}RR=D_{\rm rows}>0$.
This proves \texttt{RRRankF} on the same family.

The dual used below is the literal algebraic dual lattice
$RR^\vee=\{\lambda:\lambda(RR)\subset\mathcal O_K\}$ with its
stored dual metric.  A finite generating set and a nonzero
denominator show that this is finitely generated and spans the
full dual space.  There is no trace dual or different ideal in
this assertion.  For a saturated sublattice the determinant
identity for a lattice and its metric quotient proves
$\deg(Q^\vee)=-\deg(Q)$.  Annihilators identify saturated
subobjects with dual quotients, reversing their ranks and
slopes.  Taking the largest and smallest polygon slopes gives
\[
 \mu_{\min}(RR^\vee)=-\mu_{\max}(RR).
\]
For each full positive-rank $N_u$, its mean slope is the weighted
mean of its polygon slopes and is at most its maximum.
Together with \texttt{NuFullF} and \texttt{MuRank}, these are
exactly the two fields of \texttt{Lat\_at}; a common maximum
threshold proves \texttt{LatF}.
% END GFA6

% BEGIN GFA7
\paragraph{GFA7: putting the original maps in the common Gram calculation.}
\label{claim:GFA7}
Keep the hypotheses of GFA6, now with $C_0\ge1$, and supply
the original full and retained lattices.  The period normalization is
\[
 c_u=\frac{\operatorname{vol}(N_u^{\rm an})
                   \operatorname{vol}(H_u^{\rm an})}{\dim\Pi_u}.
\]
It is nonnegative by the source's positive volumes and dimension
formula.  If $c_u=0$, the source's normalized kernel is zero,
again contradicting its diagonal defect $e_0<1$ and $d_u>0$.
Thus $c_u>0$ for every word.  With $A_u^{\rm an}=c_u^{-1/2}P_u$,
the Gram blocks in GF1 are exactly
$(A_u^{\rm an})^*A_{u'}^{\rm an}$.

We spell out the coordinate identification needed to use GF8.
Choose an orthonormal basis of the original retained target
subspace and extend it to the full analytic target.  Its first
block has dimension $\operatorname{rank}RR$, by GFA6.
Choose orthonormal bases of every full source and every retained
source.  The inclusion matrix $V_u$ of the retained space has
$V_u^*V_u=1$.  The original \texttt{BesselId} in
\texttt{LinkDIV} identifies the projected analytic map with the
scalar extension of the original rational map
\[
 B_u:N_u\longrightarrow RR^\vee.
\]
The scalar in this identity is nonzero, as explicitly required
there; absorb it and $c_u^{1/2}$ into the invertible source
coordinate matrix.  The target coordinate change is the inverse
transpose of the primal basis change, as required by the dual
basis.  Thus, at the actual designated complex embedding,
\[
 YB_u^\sigma=P_0 A_u^{\rm an}J_uX_u,
 \qquad J_u^*J_u=1,\quad X_u,Y\text{ invertible}.             \tag{GFA.18}
\]
If the representative infinite-place embedding is the conjugate
of the designated embedding, conjugate every matrix in this
equality.  Conjugation preserves inverses, products, Gram
matrices and squared norms, so the same equality holds there.
This verifies the source's possible identity-or-conjugation
choice, without choosing unrelated embeddings for different
words.

The remaining geometric requirements in GF8 are also literal
fields of the supplied records.  \texttt{OneMapIntoCC} gives
$B_u(N_u)\subset RR^\vee$.  The nonzero-map height bound of
\texttt{LinkDIV} gives the normalized squared height
$h(B_u)\le D.\texttt{div.h}(u)$ for every nonzero $B_u$;
the zero maps impose no height condition.  The target is the
same $RR^\vee$ as in GFA6.  Its retained fraction is
$1-\delta$.  The isometries $J_u$ use the original retained
source ranks $r_u$, not the full $d_u$.  Finally
\texttt{LinkS}, activated by the same \texttt{ChosenProviso},
supplies all the mixed and diagonal Gram estimates and
$D_{\rm all}\le(1+\eta_1)D_*$.
Replacing $\eta_0,\eta_1$ by their nonnegative parts weakens
these inequalities and preserves their uniform upper bounds.
Consequently every object, normalization, and map hypothesis
in GF1--GF8 has now been matched to this one original family.
% END GFA7

% BEGIN GFA8
\paragraph{GFA8: the literal frame inequality for every height parameter.}
\label{claim:GFA8}
Let $\sigma=x^{-C_H}$, and suppose $0<\sigma\le1/20$,
$0\le d\le1/4$, and
$0\le\eta_0,\eta_1,\sqrt{e_0}\le\sigma/2$.
The proviso and code condition give
\[
 \sum_ud_u\ge\frac{D_*}{\sigma},
 \qquad \sum_u(d_u-r_u)\le d\sum_ud_u.                       \tag{GFA.19}
\]
The first follows by multiplying
$d_u\ge D_*/(C_0\Gamma^2)$ by
$\#\mathrm{Word}\ge C_0\Gamma^2x^{C_H}$.

Define the actual frame quantities
\[
 s=\frac{\sum_ur_u(\mu(N_u)-h_u)}{\sum_ur_u},\qquad
 u_* =\max_u(\mu_{\max}(N_u)-h_u)-s\ge0.                    \tag{GFA.20}
\]
The nonnegativity follows from GFA6 and positive ranks; the
denominator is positive because the word set is nonempty.
Every adjusted source polygon slope is at most $s+u_*$, and
their weighted average is exactly $s$.
For \emph{any} $H>0$ satisfying
$s-\mu_{\min}(RR^\vee)\le H$, all target slopes are at least
$s-H$.  Apply GF8 with the objects and heights in GFA7:
\[
 s-\frac{4u_*H/(u_*+H)+(5\sigma+4d+4\sqrt{e_0})H}{1-\delta}
       \le\mu(RR^\vee).                                    \tag{GFA.21}
\]
For $u_*>0$, integration of $(1+t)^{-1}$ on $[0,H/u_*]$
gives
$H/(u_*+H)\le\log(1+H/u_*)$.
If $u_*=0$, both multiplied expressions are zero; this also
agrees with the source's total real-division convention.
Since $4d\le3\sqrt d$ for $d\le1/4$ and
$\sqrt{e_0}\le\sqrt{2e_0}$, (GFA.21) implies
\[
 s-\log\#\mathrm{Word}
 -\frac{4u_*\log(1+H/u_*)+
        (5\sigma+3\sqrt d+4\sqrt{2e_0})H}{1-\delta}
 \le\mu(RR^\vee).                                         \tag{GFA.22}
\]
The extra logarithm is nonnegative because the word set is
nonempty.  This proves the exact \texttt{Frame\_at} predicate,
with witnesses $\eta=5\sigma+3\sqrt d$, $e=2e_0$.
The proof does not restrict $H$ to one later application or
assume a numerical upper bound for it.  In particular it supplies
the full universal quantifier required by that predicate.
% END GFA8

% BEGIN GFA9
\paragraph{GFA9: one pair of frame errors before the family member.}
\label{claim:GFA9}
Assume the preceding family inputs, \texttt{TuSize}, and the
actual ramified rank formula.  Write
$E(x)=\exp(-x^{\epsilon/8})$.
Bernoulli and $T_u\ge b$ give
\[
 d_u-r_u\le(2f/b)d_u.
\]
The rank windows imply $f\le x^5$; hence, eventually,
$2f/b\le E(x)^2$, uniformly for every $i,u$.
Indeed
$\log(2x^5)-x^\epsilon\le-2x^{\epsilon/8}$ eventually.
Thus in GFA8 we can choose
\[
 d=E(x)^2,\qquad e_B=2E(x)^2,\qquad
 \eta_B=5x^{-C_H}+3E(x),                                  \tag{GFA.23}
\]
before choosing $i$.  The spectral errors satisfy
$\eta_0,\eta_1,\sqrt{e_0}\le E(x)$ on every member.
For fixed $C_H>0$, enlarge a single threshold until
\[
 x^{-C_H}\le1/20,\quad E(x)\le x^{-C_H}/2,
 \quad 2E(x)^2\le1/4,
 \quad (3+4\sqrt2)E(x)x^{C_H}\le1.
\]
All follow from exponential domination of a fixed power.
Now $\eta,e$ from GFA8 obey
$0\le\eta\le\eta_B$, $0\le e\le e_B$, and
\[
 (\eta_B+4\sqrt{e_B})x^{C_H}
 =5+(3+4\sqrt2)E(x)x^{C_H}\le6.                            \tag{GFA.24}
\]
Equations (GFA.22)--(GFA.24) are the complete literal
$\texttt{FrameF}(\mathcal F,C_H)$.
Its source discard $d$ and target discard $\delta$ remain
separate: only the latter occurs in the denominator
$(1-\delta)^{-1}$.  The fixed normalization constants
$C_{\rm abs}=2$, $C_0=3$, $C_5=4$, $e_D=2$, $e_L=0$
have not been altered.
% END GFA9

% BEGIN GFA10
\paragraph{GFA10: the original thick, top, and thin nodes.}
\label{claim:GFA10}
We prove \texttt{NodesF} for any supplied original family.
Put $w''=b'/(xn)$ and $\eta=L.\mathrm{etaSz}$.
The size record gives $w''\ge\sqrt b\ge1$, so $b'\ge1$.
The parameter condition $t\eta\le w''/x$ and $t\ge1$ give
$\eta\le w''/x\le b'/3$ for $x\ge3$.
The nonnegative second summand of the original $\rho$ inequality
shows
\[
 \frac{N_*x^\epsilon\Lambda}{k}\le\frac{x^{-2}}4\le\frac1{36}.
\]
Consequently the original bottom window gives
$b_t\le(1+1/36)b'$.  It follows that
\[
 w''+\eta\le2b',\qquad
 b_t+b'/x+\eta\le\frac{61}{36}b'\le2b'.                     \tag{GFA.25}
\]

Here is the exact place dictionary in \texttt{Nodes\_at}.
For $\rho=((0,\rho_2),h)\in\Phi$, set
\[
 g(\rho)=((-1)^{\rho_2},1,\iota h)\in D.R.\Psi,
 \qquad \tau=D.eTau^{-1}(g(\rho)).
\]
Its code is exactly $\texttt{mcode}(\rho)$, so the stipulated
reading of $w(j,\rho)$ is the original $w_{j+1}(\tau)$.
The thick condition says $\rho=((0,1),t+u_0)$ with
$t\in T_2$.  The source hypotheses identify $T_2$ with the
digits of the original integer second set and identify $u_0$
by its injective integer label.  Therefore $g(\rho)$ is thick
exactly when $\rho$ is thick.  This follows by comparing the
two sign coordinates and the injective $\iota$ coordinate;
it does not follow merely from equality of cardinalities.

At a thick place the stored interval is
$b\le w_{j+1}(\tau)\le b+\eta+1$.
At a non-thick top it is
$b_t+b'/x\le w_{j+1}(\tau)\le b_t+b'/x+\eta$;
here $b_t\ge b'\ge1$ and (GFA.25) applies.
At a place that is neither thick nor top it is
$w''\le w_{j+1}(\tau)\le w''+\eta$, so again
$1\le w_{j+1}(\tau)\le2b'$.
These cases prove all three clauses of \texttt{Nodes\_at}:
all nodes in $\Phi$ are at least $1$, thick nodes have their
stated $b$-interval, and the remaining nodes are at most $2b'$.
No clause for an unrepresented code is needed on $\Phi$ because
the displayed inverse gives a representative there.  The
definition's extra convention $w=1$ elsewhere is retained.
The threshold $\max(3,\mathcal F.x_1+1)$ proves
\texttt{NodesF}.
% END GFA10

% BEGIN GFA11
\paragraph{GFA11: all seventeen original size-window clauses.}
\label{claim:GFA11}
Supply the same family's \texttt{ThickSetDensityEventually}.
For every fixed $c_6>0$ we prove \texttt{WinF} with $e_K=5$.
The raw \texttt{SizeM}, \texttt{SizeK}, \texttt{SizeN}
windows have exponents
\[
 m\in[x^{1/2-\epsilon/16-\epsilon/100},
       x^{1/2-\epsilon/16+\epsilon/100}],\quad
 k\in[x^{9/2+10\epsilon-\epsilon/100},
       x^{9/2+10\epsilon+\epsilon/100}],\quad n\le x^3.
\]
Since $0<\epsilon<1/52$, comparison of these exponents gives
\[
 x^{1/2-\epsilon}\le m\le x^{1/2+\epsilon/2},\qquad
 xm\le k\le x^5,
 \qquad k\ge x^4.                                         \tag{GFA.26}
\]
The separate parameter field $k_{\rm up}$ reads exactly
$k\le n^2x^{1/2-\epsilon}$.  This is retained as a different
inequality; it is not inferred from the loose windows in
(GFA.26).

Eventually $k\le x^5\le X'\le a$, whence $kb'\le a$.
GFA10 gives $w''\ge1$, so $X''=a/w''\le a$.
The top interval
$a\le a_t\le a+a/(kx^{13})$ proves both
$a_t\le2a$ and
$0\le a_t/b-x\le x/(kx^{13})\le1$.
The literal code-size upper bound
$\#\mathrm{Word}\le2^{k-3}$ and its nonempty word set give
$\log\#\mathrm{Word}\le(k-3)\log2$.

We also check $\mathcal L_a=L.\mathrm{La}=\log(a\,\mathrm{rd}_E)$,
which differs from $\Lambda=L.\mathrm{LaG}$.
The size record has $\mathrm{rd}_E\ge2$ and
\[
 \mathrm{rd}_E\le\frac{Q_S}{2}\sqrt{qd},\quad
 Q_S\le4n x^{\epsilon/100},\quad
 \log q,\log d\le x^{\epsilon/5+\epsilon/100}.
\]
Hence
\[
 0\le\mathcal L_a
 \le x^\epsilon+5\log x+x^{\epsilon/5+\epsilon/100}
 \le2x^\epsilon                                              \tag{GFA.27}
\]
eventually.  In the middle inequality use
$\log n+(\epsilon/100)\log x\le3\log x$ from the actual
\texttt{SizeN} upper exponent and $\log2\le\log x$.
In the last use $5\log x\le x^\epsilon/2$ and
$2x^{\epsilon/5+\epsilon/100}\le x^\epsilon$.
The same record gives $v=\mathrm{rd}_E/2>0$.

There are two small-error clauses.  Since $f=2n\le2x^3$,
choose a threshold with $x^7\le e^{c_6x^{\epsilon/8}}$.
For $x\ge3$, $4(x+1)\le x^3$, and therefore
\[
 2f(x+1)e^{-c_6x^{\epsilon/8}}\le1/x.                       \tag{GFA.28}
\]
The source's different parameter inequality
$\eta\le b/(t^3fk^3)$, with $t\ge1,k\ge x^4$, gives
$f\eta x^{12}\le b$.  Since
$12x(x+1)\le x^{12}$ for $x\ge3$,
$12x(x+1)f\eta\le b$.
Also $12x(x+1)f\le x^9\le b$ eventually.
Adding and dividing by $2x$ proves
\[
 6f(x+1)(\eta+1)\le b/x.                                  \tag{GFA.29}
\]

The seventeen fields of $\texttt{Win\_at}(L,5,c_6)$, in their
source order, are now accounted for as follows:
\[
\begin{array}{c|l}
1,2,3 & m\text{ upper},\ xm\le k,\ k\le x^5\quad\text{by (GFA.26)}\\
4,5,6 & k\le a,\ 0\le\mathcal L_a,\ \mathcal L_a\le2x^\epsilon\\
7,8 & kb'\le a,\ X''\le a\\
9,10 & a_t\le2a,\ \log\#\mathrm{Word}\le(k-3)\log2\\
11,12 & \text{(GFA.28), (GFA.29)}\\
13,14 & m\text{ lower in (GFA.26)},\quad k\le n^2x^{1/2-\epsilon}\\
15 & x^{-\epsilon}\le |T_2|/n\quad\text{by the supplied thick density}\\
16,17 & |a_t/b-x|\le1,\quad v>0.
\end{array}
\]
All exponential comparisons involve fixed positive exponents
and fixed constants.  Their finite maximum, together with the
thick-density and original family thresholds, precedes $i$.
This proves the literal \texttt{WinF}, with the prescribed
order $\forall c_6>0\ \exists e_K,x_0\ \forall x\ge x_0\ \forall i$.
% END GFA11

% BEGIN GFA12
\paragraph{GFA12: the original primitive-character reading and \texttt{CharF}.}
\label{claim:GFA12}
There is an elementary uniqueness observation useful for the
exact source choices.  Suppose primitive Dirichlet characters
$\chi_1,\chi_2$ agree at every prime outside the divisors of
a fixed nonzero integer $M$.  For any unit residue class modulo
$N_1N_2$, CRT gives a positive representative $a$ avoiding
every prime dividing $M$: retain the given unit residue at
the primes dividing $N_1N_2$, and impose $a\equiv1$ at the
other prime divisors of $M$.  Every prime factor of $a$ is
outside $M$.  Complete multiplicativity gives
$\chi_1(a)=\chi_2(a)$.  Thus their inductions to $N_1N_2$
are equal.  Factoring that character through the gcd of its
two primitive moduli proves $N_1=N_2$ and equality of the
primitive characters, hence equality at \emph{all} positive
integers.  More explicitly, CRT decomposes the common character
into characters of the groups of units modulo prime powers.
If such a local character factors through both exponents,
it factors through their minimum: the kernel at that minimum
is the larger of the two nested reduction kernels.  Taking
the product proves the asserted gcd factorization and the
uniqueness of the minimal conductor.  No prime-in-progression existence theorem is needed
for this uniqueness argument.

On the actual host, AH9--AH10 identify the specified Galois
map and Artin map, including their inertia and Frobenius
cosets.  Let
\[
 t_E:D.R.GG\longrightarrow(\mathbb Z/2\mathbb Z)^2\times D.R.A
\]
be the source's exact multiplicative-to-additive isomorphism.
It sends the second multiplicative sign to the first additive
sign and the first multiplicative sign to the second, as in
\texttt{gcode}.  For an additive character $\phi$, set
$c(\phi)(g)=\phi(t_Eg)$; the inverse is precomposition by
$t_E^{-1}$.  These maps are inverse character isomorphisms,
and $c(\phi)=1$ exactly when $\phi$ is the trivial additive
character (written $0$ in the source).

The original reciprocity map is surjective.  WB1 applied to
$c(\phi)\circ\mathrm{art}$ supplies its primitive inducer
$\chi_\phi$, with the precise values
\[
 \chi_\phi(n)=\mathrm{primVal}(\mathrm{arithOf}(D.host),\phi,n)
 \qquad(n\ge1).                                           \tag{GFA.30}
\]
Its values at unramified primes are the required Frobenius
values by AH9.  Therefore it is an actual witness of
\texttt{Reads}, proving that the source's primitive-reading
choice has a witness.  Every possible such reading agrees
with this one off the finite discriminant set.  The preceding
uniqueness proof identifies the source's stored
$\mathrm{prim}(D,c(\phi))$ with this same primitive character,
including every ramified positive integer in (GFA.30).
Surjectivity of reciprocity also gives
$\chi_\phi=1$ if and only if $\phi=0$.

Define the witness required by \texttt{Char\_at} to be
\[
 \ell(\phi)=\mathrm{Lp}(D,c(\phi)).
\]
If any primitive $\chi$ is prescribed by all the positive
integer values in (GFA.30), the same uniqueness argument
identifies it with the stored character, so
$\ell(\phi)=L'/L(1,\chi)$ for $\phi\ne0$.
The centered ramps also agree exactly:
\[
 \mathrm{GPsi}(\mathrm{arithOf}(D.host),\ell,\phi,y)
       =\mathrm{Gc}(D,c(\phi),y).                           \tag{GFA.31}
\]
For a nontrivial character substitute (GFA.30) in the finite
von Mangoldt sum and the same $y\ell(\phi)$ term.  Its
positive-part factor equals $y/n-1$ on the actual summation
range $1\le n\le\lfloor y\rfloor$.  If that range is empty,
both sums are zero.  For the trivial character both sides
are the same centered principal ramp, independently of the
unused value $\ell(0)$.  This checks the centering, including
$0<y<1$.

The distinguished additive characters $\chi_C,\psi_C$ map
under $c$ to the second and first multiplicative signs,
respectively.  Their product maps to the product of these
signs.  The original host includes the field
\texttt{Host.flPS}, which supplies its actual \texttt{PsQ}
quadratic-root identities; AH9--AH10 identify their Artin
and Frobenius readings on every such host.  These identities
identify the corresponding primitive characters with the
inducers of $\operatorname{chDir}(d)$,
$\operatorname{psDir}(q)$ and their product, in the source's
modulus-$8qd$ reading.  For $d\equiv3\pmod4$, quadratic
reciprocity identifies $\operatorname{chDir}(d)$ with the
Legendre character modulo $d$, which is the actual
$\mathrm{chiK}(D)$.  Thus the definition of $\mathrm{dlogL}$
and the stored logarithmic-derivative field give
\[
 L.\mathrm{logDerivativeAtOne}=\operatorname{Re}\ell(\chi_C).
                                                               \tag{GFA.32}
\]
Only this real-part equality is required here.

Every $\chi_\phi$ has the source's \texttt{IsCharED}
witness $c(\phi)$: at integers prime to the host modulus,
(GFA.30) equals its value on the actual Artin image.
Consequently \texttt{ChThin\_at} applied to these very
characters and (GFA.31) gives \texttt{ThinScales}, with the
same $X'=e^{x^{\epsilon/3}}$, the same $c_6$, and the same
lower limit $y\ge X'/3$.
Likewise the three clauses of \texttt{ChQD\_at}, using the
three just-identified quadratic readings, give
\texttt{ValuesAtX} with its original $C_5$.
Equations (GFA.30)--(GFA.32) and these two implications are
all four fields of \texttt{Char\_at}.  They are pointwise
identities for every original record; thus \texttt{CharF}
has threshold zero for every supplied $c_6,C_5$.
The thin and quadratic estimates remain their stated
\texttt{ChosenThin}/\texttt{ChosenQD} inputs.
% END GFA12

\paragraph{The remaining final-global work after these installations.}
On the same original family, GVC supplies \texttt{CountF} and
\texttt{SlopesMF}; GFA1--GFA5 supply \texttt{PlacesF} and
\texttt{SlopesNF}; GFA6--GFA9 supply \texttt{RRRankF},
\texttt{LatF} and the literal \texttt{FrameF}; GFA10--GFA12
supply \texttt{NodesF}, \texttt{WinF} and \texttt{CharF}.
These are proved implications from their stated original
inputs, not new producers for those inputs.

The final global obligation still has two substantive parts:
the original $\texttt{LossFS}(\mathcal F,0)$ estimate from
the chosen small shares, and the exact common-family
\texttt{linkCore}/\texttt{ChainTied} installation with those
losses.  The latter must carry the same selected top and
bottom values, original host, characters, all eleven ledger
fields, and the fixed constants through the final
\texttt{Closed} wrapper.  ZW1--ZW4 already prove the numerical
and analytic consequence once that actual tied ledger is
available.  Therefore \texttt{TASK\_GL\_FINAL\_ASSEMBLY}
remains open B with a narrower, explicit contract; it is not
declared closed by listing the now-completed inputs.


